\documentclass[10pt, oneside]{article}

% =========================
% Chinese font with ctex
% Compile with XeLaTeX
% =========================
\usepackage[UTF8, scheme=plain, fontset=none]{ctex}

\setCJKmainfont{STFangsong}     

% =========================
% Page geometry
% =========================
\usepackage[
	letterpaper,
	tmargin=1in,
	bmargin=1in,
	lmargin=1.1in,
	rmargin=1.1in
]{geometry}

% =========================
% Math
% =========================
\usepackage{amsmath, amsthm, amssymb, mathrsfs, verbatim, bbm, color, graphicx, authblk}
\usepackage{booktabs, multirow, makecell,colortbl, enumitem}
\usepackage{braket}
\usepackage{caption, subcaption}
\usepackage{framed}
\usepackage{bm}

% =========================
% Author / affiliation
% =========================
\renewcommand\Authfont{\fontsize{12pt}{14pt}\selectfont}
\renewcommand\Affilfont{\fontsize{12pt}{14pt}\selectfont\itshape}
\setlength{\affilsep}{0.25em}

% =========================
% Contents / section control
% =========================
\usepackage{tocloft}
\usepackage{titlesec}

% =========================
% Hyperref
% =========================
\usepackage{hyperref}

\hypersetup{
	colorlinks=true,
	linkcolor=blue,
	filecolor=blue,
	urlcolor=blue,
	citecolor=blue,
}

% =========================
% Page number: bottom center
% =========================
\makeatletter
\renewcommand{\ps@plain}{%
	\renewcommand{\@oddhead}{}%
	\renewcommand{\@evenhead}{}%
	\renewcommand{\@oddfoot}{\hfil{\fontsize{12pt}{14pt}\selectfont\thepage}\hfil}%
	\renewcommand{\@evenfoot}{\hfil{\fontsize{12pt}{14pt}\selectfont\thepage}\hfil}%
}
\makeatother

\pagestyle{plain}

% =========================
% Main font size
% 正文 12pt，正文行距 14pt
% =========================
\AtBeginDocument{\fontsize{12pt}{14pt}\selectfont}

% =========================
% Paragraph
% =========================
\setlength{\parindent}{2em}
\setlength{\parskip}{0pt}

% =========================
% Section / subsection labels
% section: I. II.
% subsection: A. B.
% =========================
\renewcommand{\thesection}{\Roman{section}}
\renewcommand{\thesubsection}{\Alph{subsection}}

% section 标题：左对齐
\titleformat{\section}
	{\normalfont\bfseries\fontsize{12pt}{14pt}\selectfont}
	{\thesection.}
	{1.2em}
	{}

% subsection 标题：往后缩进一点
\titleformat{\subsection}
	{\normalfont\bfseries\fontsize{12pt}{14pt}\selectfont}
	{\thesubsection.}
	{1.2em}
	{}

% section 上下间距
% {左缩进}{标题前距离}{标题后距离}
\titlespacing*{\section}
	{0pt}
	{3.0ex plus 0.5ex minus 0.2ex}
	{2.8ex plus 0.3ex}

% subsection 上下间距
% subsection 往右缩进 2em，并且标题后留出距离
\titlespacing*{\subsection}
	{2em}
	{3.0ex plus 0.5ex minus 0.2ex}
	{2.2ex plus 0.3ex}

% =========================
% Contents style
% =========================
\renewcommand{\contentsname}{CONTENTS}

% CONTENTS 标题
\renewcommand{\cfttoctitlefont}{\hspace*{1.4em}\fontsize{11pt}{12pt}\selectfont\bfseries}
\renewcommand{\cftaftertoctitle}{}

\setlength{\cftbeforetoctitleskip}{0pt}
\setlength{\cftaftertoctitleskip}{0.65em}

% 目录无点线
\renewcommand{\cftdotsep}{\cftnodots}

% section / subsection 目录缩进
\setlength{\cftsecindent}{0pt}
\setlength{\cftsecnumwidth}{2.0em}

\setlength{\cftsubsecindent}{2.0em}
\setlength{\cftsubsecnumwidth}{2.0em}

% 目录编号格式：
\renewcommand{\cftsecaftersnum}{.}
\renewcommand{\cftsubsecaftersnum}{.}

% 目录字体
\renewcommand{\cftsecfont}{\fontsize{14pt}{16pt}\selectfont}
\renewcommand{\cftsecpagefont}{\fontsize{14pt}{16pt}\selectfont}
\renewcommand{\cftsubsecfont}{\fontsize{14pt}{16pt}\selectfont}
\renewcommand{\cftsubsecpagefont}{\fontsize{14pt}{16pt}\selectfont}

\setlength{\cftbeforesecskip}{0.8em}
\setlength{\cftbeforesubsecskip}{0.2em}

% \numberwithin{equation}{subsection}

% =========================
% Title information
% =========================
\title{\fontsize{15pt}{20pt}\selectfont\bfseries
转角石墨烯的连续模型}

\author[1]{Qian-Sheng Shi}

\affil[1]{
School of Physical Science and Technology, ShanghaiTech University, \protect\\
Shanghai 201210, China.
}

\newcommand{\notestartdate}{July 11, 2026}

\date{\fontsize{11pt}{14pt}\selectfont
(Dated: \notestartdate --\today)}

% =========================
% Custom maketitle style
% 不让 \maketitle 自动换页
% =========================
\makeatletter

\renewcommand{\@maketitle}{%
	\begin{center}
		{\@title \par}
		\vspace{1em}
		{\@author \par}
		\vspace{0.15em}
		{\@date \par}
	\end{center}
}

\renewcommand{\maketitle}{%
	\begingroup
		\let\footnote\thanks
		\@maketitle
	\endgroup
}

\makeatother

\newcommand{\te}{\theta}
\newcommand{\ph}{\phi}
\newcommand{\ps}{\psi}
\newcommand{\vph}{\varphi}
\newcommand{\vps}{\varpsi}
\newcommand{\pr}{\prime}
\newcommand{\pa}{\partial}
\newcommand{\mr}[1]{\mathrm{#1}}
\newcommand{\mb}[1]{\mathbf{#1}}
\newcommand{\p}[1]{\left(#1\right)}
\newcommand{\bb}[1]{\left[#1\right]}
\newcommand{\br}[1]{\left\{#1\right\}}
\newcommand{\bv}[1]{\left|#1\right|}

% =========================
% Document
% =========================
\begin{document}

{\fontsize{12pt}{16pt}\selectfont
\tableofcontents
}

\vspace{0.8em}

\maketitle

\thispagestyle{plain}

\clearpage

这里我们给出一个完全从连续模型的定义和对称性的角度出发来推导BM模型，同时理清BM模型中各种符号规范。

我们做如下规范：在本文中，统一将\(d\)维空间中的积分测度记为\(d \bm{r} \equiv d^{d} \bm{r}\)。在二维系统当中， \(d \bm{r} \equiv d^{2} \bm{r} = d x dy\)

同时，对于任意矢量\(\bm{a} = \left(a_i, a_j, a_k\right)\)
, 对于矢量其本身采用粗体，而对于矢量在各个方向上的分量不采用粗体进行标记。
\section{连续模型范式}

我们首先可以通过 Fourier 变换构造一个定义在连续空间上的插值场，使其在格点处重现格点算符。进一步投影到某个低能动量区域后，该插值场成为缓慢变化的包络场；在这一低能和长波条件下，格点求和可以近似为连续积分，并对包络场进行梯度展开，从而得到连续低能模型。

紧束缚模型的哈密顿量通常表示为
\begin{align}
    \hat{H} = \sum_{\bm{R},\bm{R}^{\prime}} t \left(\bm{R}- \bm{R}^{\prime}\right) \hat{c}^{\dagger}_{\bm{R}} \hat{c}_{\bm{R}^{\prime}},
\end{align}
其中\(\bm{R}\)是格矢，可以被表示成\(\bm{R} = n_1 \bm{a}_{1} + n_2 \bm{a}_{2}\)。这里的格点的费米子算符满足反对易关系
\begin{align}
	\left\{\hat{c}_{\bm{R}}, \hat{c}^{\dagger}_{\bm{R}^{\prime}}\right\} = \delta_{\bm{R}, \bm{R}^{\prime}}.
\end{align}
定义 \(\bm{\delta} = \bm{R} - \bm{R}^{\pr}\), 则有\(\bm{R}^{\prime} = \bm{R} - \bm{\delta}\)。 所以哈密顿量可以重新写成
\begin{align}
	\hat{H} = \sum_{\bm{R}, \bm{\delta}} t \left(\bm{\delta}\right) \hat{c}^{\dagger}_{\bm{R}} \hat{c}_{\bm{R} - \bm{\delta}}. 
\end{align} 
现在我们需要将离散的格点算符转换成连续场算符，之所以需要连续场算符\(\hat{\psi}\), 是因为我们需要用一个定义在连续空间上的函数，来对离散的格点进行插值。格点的费米子算符的傅里叶变换为
\begin{align}
	\hat{c}_{\bm{R}} = \frac{1}{\sqrt{N}} \sum_{\bm{k} \in \mathrm{BZ}} e^{i \bm{k} \cdot \bm{R}} \hat{c}_{\bm{k}}.
\end{align}
现在定义一个作用在连续坐标\(\bm{r}\)上的插值场，
\begin{align}
	\hat{\Psi} \left(\bm{r}\right) = \frac{1}{\sqrt{\mathcal{A}}} \sum_{\bm{k} \in \mathrm{BZ}} e^{i \bm{k} \cdot \bm{r}} \hat{c}_{\bm{k}},
\end{align}
其中\(\mathcal{A}\)是体系的面积（这里仅讨论二维情况）。这里\(\bm{r}\)可以取任意连续值，而并不要求是格点，
当取
\begin{align}
	\bm{r} = \bm{R},
\end{align}
时， 
\begin{align}
	\hat{\Psi} \left(\bm{R}\right) = \frac{1}{\sqrt{\mathcal{A}}} \sum_{\bm{k} \in \mathrm{BZ}} e^{i \bm{k} \cdot \bm{R}} \hat{c}_{\bm{k}},
\end{align}
对于周期体系的面积，有
\begin{align}
	\mathcal{A} = N A_{\mathrm{uc}},
\end{align}
这里\(N\)是原胞数量，\(A_{\mathrm{uc}}\)是单个原胞的面积。于是格点的费米子算符可以表示成连续函数在格点处的取值。
\begin{align}
	\boxed{
		\hat{c}_{\bm{R}} = \left(A_{\mathrm{uc}}\right)^{1/2} \hat{\Psi} \left(\bm{r}\right) \bigg|_{\bm{r} = \bm{R}}. 
	}
\end{align}
但是这里的连续场并不一定是慢变场。格点费米子算符需要用连续包络场在格点处的取值表示，而格点求和则在低能近似(慢变近似)下改写为连续积分。离散跃迁位移\(\bm{\delta}\)仍然保留，并在后续梯度展开中形成连续哈密顿量的各阶系数。而上式中对动量\(\bm{k}\)是在整个BZ中进行求和，因此如果存在很大的动量分量，那么连续场可能在晶格尺度上快速变化，并不一定适合做梯度展开。连续插值场的存在只是傅里叶变换的结果。

现在引入低能动量截断，假设我们关注的低能态在某个动量\(\bm{K}\)附近，我们可以将原来位于布里渊区中的动量写成\(\bm{K}\)和一个小动量\(\bm{q}\)的求和，也就是
\begin{align}
	\bm{k} = \bm{K} + \bm{q},
\end{align}
其中我们对小动量的区域做一定的限制，取截断\(\Lambda\) 
\begin{align}
	\left|\bm{q}\right| < \Lambda, \quad \Lambda a \ll 1,
\end{align}
\(a\) 是周期性体系的晶格常数。在这个小区域内，格点算符可以近似成
\begin{align}
	\hat{c}_{\bm{R}}  \simeq \frac{1}{\sqrt{N}} \sum_{\left|\bm{q}\right|  < \Lambda} e^{i \left(\bm{K} + \bm{q} \right) \cdot \bm{R}} \hat{c}_{\bm{K} + \bm{q}}.
\end{align}
\(e^{i \bm{K} \cdot \bm{R}}\) 是一个实空间中快速振荡的相位因子（对于本文后续关心的非零低能动量中心，我们通常假设\(\bm{K} \sim a^{-1}\)），提出该相位，则有
\begin{align}
	\hat{c}_{\bm{R}}  \simeq \frac{1}{\sqrt{N}}   e^{i \bm{K} \cdot \bm{R}} \sum_{\left|\bm{q}\right|  < \Lambda}  e^{i \bm{q} \cdot \bm{R}} \hat{c}_{\bm{K} + \bm{q}},
\end{align}
定义低能慢变包络场
\begin{align}
	\boxed{
		\hat{\psi}\left(\bm{r}\right) = \frac{1}{\sqrt{\mathcal{A}}} \sum_{\left|\bm{q}\right|  < \Lambda}  e^{i \bm{q} \cdot \bm{r}} \hat{c}_{\bm{K} + \bm{q}},
	}
\end{align}
可以看到包络场的空间变化由\(\bm{q}\)控制。由于所选取的低能动量中心满足\(\left|\bm{K}\right| \sim a^{-1}\), 而包络动量\(\left|\bm{q}\right| \ll a^{-1}\), 载波相位\(e^{i \bm{K} \cdot \bm{R}}\)在相邻格点之间可以产生\(\mathcal{O}\left(1\right)\)的相位变化，因此在晶格尺度上快速振荡；相比之下，\(e^{i \bm{q} \cdot \bm{R}}\)在相邻格点之间的相位变化满足\( \left|\bm{q}\right| a \ll 1\), 从而构成缓慢变化的包络场。这也是后面我们将求和写成积分，以及梯度展开的基础。

格点费米子算符变成
\begin{align}
	\boxed{
		\hat{c}_{\bm{R}} \simeq \left(A_{\mathrm{uc}}\right)^{1/2}  e^{i \bm{K} \cdot \bm{R}} \hat{\psi}\left(\bm{R}\right)
	}
\end{align}
带入到哈密顿量中有
\begin{align}
	\hat{H} = A_{\mathrm{uc}} \sum_{\bm{R}, \bm{\delta}} t \left(\bm{\delta}\right) e^{ - i \bm{K} \cdot \bm{\delta}} \hat{\psi}^{\dagger} \left(\bm{R}\right) \hat{\psi} \left(\bm{R} - \bm{\delta}\right),
\end{align}
利用前面的积分求和关系\(A_{\mathrm{uc}} \sum_{\bm{R}} \rightarrow \int d \bm{r}\), 由于\(\hat{\psi}\left(\bm{r}\right)\)是慢变的，其在一个原胞尺度内变化很小，所以在一个原胞内可以认为
\begin{align}
	A_{\mathrm{uc}} \sum_{\bm{R}} F \left(\bm{r}\right) \bigg|_{\bm{r} = \bm{R}} \simeq \int d \bm{r} F \left(\bm{r}\right),
\end{align}
其中定义
\begin{align}
	F \left(\bm{r}\right) = \sum_{\bm{\delta}} t \left(\bm{\delta}\right) e^{ - i \bm{K} \cdot \bm{\delta}} \hat{\psi}^{\dagger} \left(\bm{R}\right) \hat{\psi} \left(\bm{R} - \bm{\delta}\right).
\end{align}
更严格的写法是引入Poisson求和公式（Dirac comb formula），也就是
\begin{align}
	\boxed{
		\sum_{\bm{R}}  \delta^{\left(2\right)} \left(\bm{r} - \bm{R}\right) = \frac{1}{A_{\mathrm{uc}}} \sum_{G} e^{i \bm{G} \cdot \bm{r}}
	},
\end{align}
其中\(\bm{G}\)是倒格矢，于是
\begin{align}
	A_{\mathrm{uc}} \sum_{\bm{R}} F \left(\bm{R}\right) = A_{\mathrm{uc}} \int d \bm{r} F \left(\bm{r}\right) \sum_{\bm{R}} \delta^{\left(2\right)} \left(\bm{r} - \bm{R}\right) = \int d \bm{r} F \left(\bm{r}\right) \sum_{\bm{G}} e^{i \bm{G} \cdot \bm{r}}.
\end{align}
即
\begin{align}
	A_{\mathrm{uc}} \sum_{\bm{R}} F \left(\bm{R}\right) = \int d \bm{r} F \left(\bm{r}\right) + \sum_{\bm{G} \neq 0} \int d \bm{r} e^{i \bm{G} \cdot \bm{r}} F \left(\bm{r}\right).
\end{align}
由于\(\left|\bm{G}\right| \sim 1/a\), 而\(F \left(\bm{r}\right)\)只在远长于晶格常数的尺度上变化，所以\(e^{i \bm{G} \cdot \bm{r}}\) 会在原子尺度上快速振荡，积分中正负贡献互相抵消，所以我们可以只保留\(\bm{G} = 0\)这一项。所以求和变积分，本质上是忽略晶格Poisson求和公式中的非零谐波。

那么在这个基础上，什么叫作“\(F \left(\bm{r}\right)\)是缓慢变化的呢”， 更仔细的来说，由于 \(\left|\bm{q}\right|  < \Lambda\)，\(F \left(\bm{r}\right)\)中的e指数类似为\(e^{i \left(\bm{q}^{\prime} - \bm{q}\right) \cdot \bm{r}}\), 而\(\bm{q}^{\prime} - \bm{q}\)大小最多为\(2 \Lambda\)。 只要\(2 \Lambda \ll \left|\bm{G}_{\mathrm{min}}\right|\)，意味着对所有的倒格矢\(\bm{G}\)来说，无法满足\(\bm{G} - \left(\bm{q}^{\prime} - \bm{q}\right) = 0\), 那么\(F \left(\bm{r}\right)\) 中没有接近非零原子倒格矢的傅里叶分量，所以\(\bm{G} \neq 0\)的项消失或者被强烈抑制。

综上，连续模型的哈密顿量可以写成
\begin{align}
	\boxed{
		\hat{H} \simeq \int d \bm{r} \sum_{ \bm{\delta}}  t \left(\bm{\delta}\right) e^{- i \bm{K} \cdot \bm{\delta}} \hat{\psi}^{\dagger} \left(\bm{r}\right) \hat{\psi} \left(\bm{r} - \bm{\delta}\right),
	}
\end{align}

那么接下来我们的目标就是对\(\hat{\psi} \left(\bm{r} - \bm{\delta}\right)\)做梯度展开。在满足\(\Lambda \left|\bm{\delta}\right| \ll 1\)时，可以将\(\delta\)认为是一小量，做梯度展开
\begin{align}
	\hat{\psi} \left(\bm{r} - \bm{\delta}\right) = e^{- \bm{\delta} \cdot \nabla} \hat{\psi} \left(\bm{r}\right) =  \hat{\psi} \left(\bm{r}\right) - \delta_{i} \partial_{i} \hat{\psi} \left(\bm{r}\right) + \frac12 \delta_{i}  \delta_{j}  \partial_{i} \partial_{j} \hat{\psi} \left(\bm{r}\right) + \cdots,
\end{align}
这里用到了爱因斯坦求和约定，上面的展开等价于
\begin{align}
	\hat{\psi} \left(\bm{r} - \bm{\delta}\right) = \hat{\psi} \left(\bm{r}\right) - \bm{\delta} \cdot \nabla  \hat{\psi} \left(\bm{r}\right) + \frac12 \left(\bm{\delta} \cdot \nabla\right)^{2}  \hat{\psi} \left(\bm{r}\right) + \cdots,
\end{align}
二阶项中的\(i,j\)都是重复指标，因此分别求和就是
\begin{align}
	\delta_{i}  \delta_{j}  \partial_{i} \partial_{j} = \sum_{i,j = x,y} \delta_{i}  \delta_{j}  \partial_{i} \partial_{j}  = \delta^{2}_{x} \partial^{2}_{x} + 2 \delta_{x} \delta_{y} \partial_{x} \partial_{y} + \delta^{2}_{y} \partial^{2}_{y.}
\end{align}

代回到哈密顿量
\begin{align}
	\boxed{
			\hat{H} \simeq \int d \bm{r} \hat{\psi}^{\dagger} \left(\bm{r}\right) \sum_{\bm{\delta}} t \left(\bm{\delta}\right) e^{-i \bm{K} \cdot \bm{\delta}} \left[1 -  \delta_{i} \partial_{i} +  \frac12 \delta_{i}  \delta_{j}  \partial_{i} \partial_{j} + \cdots \right] \hat{\psi} \left(\bm{r}\right) ,
	}
\end{align}
定义零阶矩
\begin{align}
	\mathcal{M}^{\left(0\right)} \left(\bm{K}\right) = \sum_{\bm{\delta}} t \left(\bm{\delta}\right)  e^{-i \bm{K} \cdot \bm{\delta}} = \varepsilon \left(\bm{K}\right).
\end{align}
即零阶矩是低能中心\(\bm{K}\)处的能量项常数。哈密顿量的零阶项为
\begin{align}
	\hat{H}^{\left(0\right)} = \int d \bm{r} \hat{\psi}^{\dagger} \left(\bm{r}\right) \varepsilon \left(\bm{K}\right) \hat{\psi} \left(\bm{r}\right) ,
\end{align}


同理，定义一阶矩
\begin{align}
	\mathcal{M}^{\left(1\right)}_{i} \left(\bm{K}\right) = \sum_{\bm{\delta}} t \left(\bm{\delta}\right)  e^{-i \bm{K} \cdot \bm{\delta}} \delta_{i}.
\end{align}
对应到实空间中一阶项是
\begin{align}
	\hat{H}^{\left(1\right)} = - \int d \bm{r} \hat{\psi}^{\dagger} \left(\bm{r}\right) \mathcal{M}^{\left(1\right)}_{i} \left(\bm{K}\right) \partial_{i}  \hat{\psi} \left(\bm{r}\right) .
\end{align}
现在我们来看一阶矩和零阶矩的关系，将零阶矩对\(\bm{k}\)求导，
\begin{align}
	\frac{\partial \varepsilon \left(\bm{k}\right)}{\partial k_{i}} = -i \sum_{\bm{\delta}} t \left(\bm{\delta}\right) e^{-i \bm{k} \cdot \bm{\delta}} \delta_{i},
\end{align}
因此在\(\bm{k} = \bm{K}\)处，
\begin{align}
	\frac{\partial \varepsilon \left(\bm{k}\right)}{\partial k_{i}}  \bigg|_{\bm{k} = \bm{K}} = - i \mathcal{M}^{\left(1\right)}_{i} \left(\bm{K}\right),
\end{align}
定义群速度
\begin{align}
	v_i \left(\bm{K}\right) = \frac{1}{\hbar} \frac{\partial \varepsilon \left(\bm{k}\right)}{\partial k_{i}}  \bigg|_{\bm{k} = \bm{K}} ,
\end{align}
所以
\begin{align}
	v_i \left(\bm{K}\right) = - \frac{i}{\hbar} \mathcal{M}^{\left(1\right)}_{i} \left(\bm{K}\right), \quad \mathcal{M}^{\left(1\right)}_{i} \left(\bm{K}\right) = i \hbar v_i \left(\bm{K}\right),
\end{align} 
因此哈密顿量的一阶项为
\begin{align}
	\hat{H}^{\left(1\right)} = \int d \bm{r} \hat{\psi}^{\dagger} \left(\bm{r}\right) \left[- i \hbar v_i \left(\bm{K}\right) \partial_{i} \right]  \hat{\psi} \left(\bm{r}\right) .
\end{align}
如果\(\bm{K}\)是费米面上所选取的低能点，可以将\(\bm{v} \left(\bm{K}\right)\)称为费米速度
\begin{align}
	\bm{v}_F = \frac{1}{\hbar} \nabla_{\bm{k}} \varepsilon \left(\bm{k}\right) \bigg|_{\bm{k} = \bm{K}}.
\end{align}
定义二阶矩
\begin{align}
	\mathcal{M}^{\left(2\right)}_{ij} \left(\bm{K}\right) = \sum_{\bm{\delta}} t \left(\bm{\delta}\right)  e^{-i \bm{K} \cdot \bm{\delta}} \delta_{i} \delta_{j}
\end{align}
对应哈密顿量的二阶项为
\begin{align}
	\hat{H}^{\left(2\right)} = \frac{1}{2} \int d \bm{r} \hat{\psi}^{\dagger} \left(\bm{r}\right) \mathcal{M}^{\left(2\right)}_{ij} \left(\bm{K}\right) \partial_{i} \partial_{j} \hat{\psi} \left(\bm{r}\right),
\end{align}
另一方面，二阶矩和零阶矩有如下关系
\begin{align}
	\frac{\partial^{2} \varepsilon \left(\bm{k}\right)}{\partial k_i \partial k_j} = - \sum_{\bm{\delta}} t \left(\bm{\delta}\right)  e^{-i \bm{k} \cdot \bm{\delta}} \delta_{i} \delta_{j},
\end{align}
所以
\begin{align}
	\frac{\partial^{2} \varepsilon \left(\bm{k}\right)}{\partial k_i \partial k_j} \bigg|_{\bm{k} = \bm{K}} = - \mathcal{M}^{\left(2\right)}_{ij} \left(\bm{K}\right).
\end{align}
定义有效质量张量
\begin{align}
	\left(\frac{1}{m}\right)_{ij} = \frac{1}{\hbar^{2}} \frac{\partial^{2} \varepsilon \left(\bm{k}\right)}{\partial k_i \partial k_j} \bigg|_{\bm{k} = \bm{K}} = -  \frac{1}{\hbar^{2}} \mathcal{M}^{\left(2\right)}_{ij} \left(\bm{K}\right).
\end{align}
即
\begin{align}
	\mathcal{M}^{\left(2\right)}_{ij} \left(\bm{K}\right) = - \hbar^{2} \left(m^{-1}\right)_{ij}.
\end{align}
因此二阶项为
\begin{align}
	\hat{H}^{\left(2\right)} = - \frac{\hbar^{2}}{2} \int d \bm{r} \hat{\psi}^{\dagger} \left(\bm{r}\right) \left(m^{-1}\right)_{ij} \partial_{i} \partial_{j} \hat{\psi} \left(\bm{r}\right),
\end{align}
如果体系在\(\bm{K}\)附近各向同性， 则
\begin{align}
	\left(m^{-1}\right)_{ij} = \frac{1}{m^{*}} \delta_{ij},
\end{align}
从而
\begin{align}
	\hat{H}^{\left(2\right)} = - \frac{\hbar^{2}}{ 2 m^{*}} \int d \bm{r} \hat{\psi}^{\dagger} \left(\bm{r}\right) \nabla^{2} \hat{\psi} \left(\bm{r}\right).
\end{align}
将零阶项，一阶项和二阶项合并得到
\begin{align}
	\hat{H}_{\mathrm{c}} \simeq \int d \bm{r}  \hat{\psi}^{\dagger} \left(\bm{r}\right) \left[\varepsilon \left(\bm{K}\right) - i \hbar v_i \left(\bm{K}\right) \partial_{i} - \frac{\hbar^{2}}{2} \left(m^{-1}\right)_{ij} \partial_{i} \partial_{j} + \cdots \right] \hat{\psi} \left(\bm{r}\right),
\end{align}
采用矢量记号，一阶项可以写为\(- i \hbar \bm{v} \left(\bm{K}\right) \cdot \nabla\)。在各向同性情况下
\begin{align}
	\boxed{
		\hat{H}_{\mathrm{c}} \simeq \int d \bm{r}  \hat{\psi}^{\dagger} \left(\bm{r}\right) \left[\varepsilon \left(\bm{K}\right) - i \hbar \bm{v} \left(\bm{K}\right) \cdot \nabla - \frac{\hbar^{2}}{2 m^{*}} \nabla^{2} + \cdots \right] \hat{\psi} \left(\bm{r}\right)
	},
\end{align}

现在将实空间中的连续模型写到动量空间，包络场的傅里叶变换为
\begin{align}
	\hat{\psi} \left(\bm{r}\right)	= \int_{\left|\bm{q}\right| < \Lambda} \frac{d \bm{q}}{\left(2 \pi\right)^{2}}  e^{i \bm{q} \cdot \bm{r}} \hat{\psi} \left(\bm{q}\right).
\end{align}
同样约定动量空间中的测度为\(d \bm{q} \equiv d^{2} \bm{q}\)。对于上式求偏导有\(- i \partial_{i} e^{i \bm{q} \cdot \bm{r}}  = q_i e^{i \bm{q} \cdot \bm{r}}\)。所以傅里叶变换在动量空间中有\(-i \partial_{i} \rightarrow q_i\)以及\(- \partial_i \partial_j \rightarrow q_i q_j\)。因此，动量空间中的连续哈密顿量可以写成
\begin{align}
	\boxed{
		\hat{H}_{c} \simeq \int_{\left|\bm{q}\right| < \Lambda} \frac{d \bm{q}}{\left(2 \pi\right)^{2}}  \hat{\psi}^{\dagger} \left(\bm{q}\right) \left[\varepsilon \left(\bm{K}\right) + \hbar v_i \left(\bm{K}\right) q_i + \frac{\hbar^2}{2} \left(m^{-1}\right)_{ij} q_i q_j + \cdots \right] \hat{\psi} \left(\bm{q}\right)
	}.
\end{align} 
因此，低能色散关系为
\begin{align}
	E \left(\bm{q}\right) = \varepsilon \left(\bm{K}\right) + \hbar v_i \left(\bm{K}\right) q_i + \frac{\hbar^2}{2} \left(m^{-1}\right)_{ij} q_i q_j + \cdots
\end{align}
注意，这里的\(\bm{q}\)是相对于低能展开中心\(\bm{K}\)的小动量，完整的晶体动量是
\begin{align}
	\boxed{
		\bm{k} = \bm{K} + \bm{q}	
	}.
\end{align}
因此可以写成
\begin{align}
	\boxed{
		E \left(\bm{q}\right) = \varepsilon \left(\bm{K} + \bm{q}\right)
	}.
\end{align}

原始紧束缚模型给出的色散关系是
\begin{align}
	\varepsilon \left(\bm{k}\right) = \sum_{\bm{\delta}} t \left(\bm{\delta}\right) e^{-i \bm{k} \cdot \bm{\delta}}.
\end{align}
注意，这里的\(\bm{\delta}\)是两个格点之间的坐标差。将\(\bm{k} = \bm{K} + \bm{q}\)代入，得到
\begin{align}
	\varepsilon \left(\bm{K} + \bm{q}\right) = \sum_{\bm{\delta}} t \left(\bm{\delta}\right) e^{-i \bm{K} \cdot \bm{\delta}} e^{-i \bm{q} \cdot \bm{\delta}}
\end{align}
将\(e^{-i \bm{q} \cdot \bm{\delta}}\)展开得到
\begin{align}
	 e^{-i \bm{q} \cdot \bm{\delta}} = 1 - i q_i \delta_{i} - \frac{1}{2} q_i q_j \delta_i \delta_j + \cdots,
\end{align}
代入到色散关系中，便可得到
\begin{align}
	\varepsilon \left(\bm{K + \bm{q}}\right) = \varepsilon \left(\bm{K}\right) + \hbar v_i q_i + \frac{\hbar^{2}}{2} \left(m^{-1}\right)_{ij} q_i q_j + \cdots .
\end{align}
因此，在（\( \Lambda a \ll 1\)）的长波条件和保留相同的展开阶数下，实空间中对包络场的梯度展开，与动量空间中围绕低能动量中心对布洛赫哈密顿量所做的长波展开模型是逐项等价的。对于多分量体系，后者即构成通常所谓的\(\bm{k} \cdot \bm{p}\)有效模型。

所以低能模型的有效性来自于两类截断，首先只保留低能动量中心\(\bm{K}\)附近满足\(\left|\bm{q}\right| < \Lambda\)的自由度；其次，在（\( \Lambda a \ll 1\)）的长波条件下，只保留梯度或者小动量展开的有限项。

 \section{转角双层石墨烯的连续模型}
按照上面连续模型的范式，这里我们给出转角双层石墨烯连续模型的推导。与前面不同的是，由于上下两层石墨烯不再共用同一套Bravais格子，我们需要用原子位置上的费米子算符写出紧束缚哈密顿量，再从原子位置上的费米子算符定义低能连续场算符，进而写出实空间的连续哈密顿量；在确认哈密顿量具有的莫尔周期性后，我们将哈密顿量在平面波基底下写出其矩阵表示。

\subsection{转角双层石墨烯的紧束缚哈密顿量}

这里我们再次重新声明三种不同的面积，首先是对应上面公式中的原胞面积\(A_{\mathrm{uc}}\), 在下文中等价于单层石墨烯原胞面积\(A_{\mathrm{mlg}}\), 同时定义莫尔超晶格原胞面积\(A_{\mathrm{M}}\)， 而\(\mathcal{A}\)为单层样品的总面积。有如下关系
\begin{align}
	\mathcal{A} = N_{0} A_{\mathrm{mlg}} = N_{\mathrm{M}} A_{\mathrm{M}}.
\end{align}
其中\( N_{0}\) 是每一层包含的石墨烯原胞数，\(N_{\mathrm{M}} \) 是莫尔原胞数。

对于未旋转时的单层石墨烯，格点位置为
\begin{align}
	\bm{R} = n_1 \bm{a}_{1} + n_{2} \bm{a}_{2},
\end{align}
其中\(n_{1},n_{2} \in \mathbb{Z}\), \( \bm{a}_{1},  \bm{a}_{2}\)是单层石墨烯的基矢，具体为
\begin{align}
	\bm{a}_{1} = a \left(\frac{1}{2}, \frac{\sqrt{3}}{2}\right),\quad \bm{a}_{2} = a \left(-\frac{1}{2}, \frac{\sqrt{3}}{2}\right),
\end{align}
这里\(a = 2.46\)\AA 为石墨烯的晶格常数。原胞内部两个原子位置为
\begin{align}
	\bm{\tau}_{A} = \left(0,0\right),\quad \bm{\tau}_{B} = a \left(0, \frac{\sqrt{3}}{3}\right) = \frac{1}{3} \left(\bm{a}_{1} + \bm{a}_{2}\right).
\end{align}
因此，对于单层石墨烯的原子位置，我们一般记为
\begin{align}
	\bm{r}_{S} = \bm{R} + \bm{\tau}_{S}.
\end{align}
其中\(S \in \left\{A,B\right\}\)是子晶格指标。
\vspace{2em}


对于转角双层石墨烯，我们引入层指标\(j \in \left\{t,b\right\}\), 记第\(j\)层的旋转角度为\(\theta_{j}\), 记上下两层的旋转角度分别为
\begin{align}
	\theta_{b} = - \frac{\theta}{2}, \quad \theta_{t} = \frac{\theta}{2}.
\end{align}
这里采用主动旋转规范，也就是正号表示逆时针旋转，负号表示顺时针旋转。所以实际每一层的原子位置写为
\begin{align}
	\boxed{
		\bm{r}_{j,S} \left(\bm{R}\right) = \mathcal{R} \left(\theta_{j}\right) \left(\bm{R} + \bm{\tau}_{S}\right),
	}
\end{align}
其中\(\mathcal{R} \left(\theta_{j}\right) = \begin{pmatrix}
	\cos \theta_{j} & - \sin \theta_{j} \\ \sin  \theta_{j}  & \cos  \theta_{j} 
\end{pmatrix}\)是逆时针旋转矩阵。

紧束缚哈密顿量应写成
\begin{align}
	\hat{H}_{\mathrm{tb}} = \sum_{j,j^{\prime}} \sum_{S,S^{\prime}} \sum_{\bm{r}_{j,S},\bm{r}^{\prime}_{j^{\prime, S^{\prime}}}} t_{j j^{\prime}} \left(\bm{r}_{j,S} -\bm{r}^{\prime}_{j^{\prime, S^{\prime}}} \right) \hat{c}^{\dagger}_{j,S,\bm{r}_{j,S} } \hat{c}_{j^{\prime}, S^{\prime}, \bm{r}^{\prime}_{j^{\prime, S^{\prime}}} }.
\end{align}
其中\( \hat{c}_{j,S,\bm{r}_{j,S} }\)是原子位置上的费米子算符，其满足反对易关系
\begin{align}
	\left\{\bm{r}_{j,S}, \bm{r}^{\prime}_{j^{\prime, S^{\prime}}} \right\} = \delta_{j,j^{\prime}} \delta_{S,S^{\prime}} \delta_{\bm{r}_{j,S}, \bm{r}^{\prime}_{j^{\prime, S^{\prime}}}}.
\end{align}
对于跃迁项，更仔细地可以写成
\begin{align}
	t_{jj^{\prime}} \left(\bm{d}\right) = t \left[\bm{d} + \left(z_{j} - z_{j^{\prime}}\right) \hat{z} \right],
\end{align}
其中 
\begin{align}
	\bm{d} = \bm{r}_{j,S}-  \bm{r}^{\prime}_{j^{\prime, S^{\prime}}}.
\end{align}
当层指标相同时，哈密顿量描述的是intralayer coupling, 当层指标不同的时候，哈密顿量描述的是interlayer tunneling. 厄米性要求
\begin{align}
	t \left(\bm{d}\right) = t^{*} \left(-\bm{d}\right).
\end{align}

\subsection{连续场算符与有效哈密顿量}

在转角双层石墨烯当中，低能展开中心也就是谷（valley），谷的位置为
\begin{align}
	\bm{K}_{j,\xi} = \mathcal{R}\left(\theta_{j}\right) \bm{K}_{\xi},
\end{align}
\(\bm{K}_{\xi}\)是未旋转的单层石墨烯的谷位置，\(\xi = \pm\) 是谷指标。对于我们所选取的convention，有
\begin{align}
	\bm{K}_{+} = \frac{4 \pi}{3 a} \left(1, 0\right), \quad \bm{K}_{-} = \frac{4 \pi}{3 a} \left(-1, 0\right),
\end{align}
在每个谷附近取阶段
\begin{align}
	\bm{p} = \bm{K}_{j, \xi} + \bm{k}, \quad \left|\bm{q}\right| < \Lambda, \quad \Lambda a \ll 1,
\end{align}
由此，定义连续场
\begin{align}
	\boxed{
		\hat{\psi}_{j,S,\xi} \left(\bm{r}\right) = \frac{1}{\sqrt{\mathcal{A}}} \sum_{\left|\bm{k}\right| < \Lambda} e^{i \bm{k} \cdot \bm{r}} \hat{c}_{j,S, \bm{K}_{j, \xi} + \bm{k}}.
	}
\end{align}
又 \(\mathcal{A} = N_{0} A_{\mathrm{mlg}}\), 原子位置费米子算符的低能展开为
\begin{align}
	\boxed{
		\hat{c}_{j,S, \bm{r}_{j,S}} \simeq \sqrt{A_{\mathrm{mlg}}} \sum_{\xi = \pm} e^{i \bm{K}_{j, \xi} \cdot \bm{r}_{j,S}} \hat{\psi}_{j,S, \xi} \left(r\right) \bigg|_{\bm{r} = \bm{r}_{j,S}}.
	}
\end{align}
(这里的近似来源于仅保留两个 valley 附近的低能 Fourier 模式，而不是来源于后续使用的 Dirac comb。对于固定子晶格的原子位置集合，Dirac comb 在保留全部原子倒格矢时是严格恒等式。) 得到有效哈密顿量为
\begin{align}
	\hat{H}_{\mathrm{eff}} = A_{\mathrm{mlg}} \sum_{j,j^{\prime}} \sum_{S,S^{\prime}} \sum_{\xi,\xi^{\prime}} \sum_{\bm{r}_{j,S}, \bm{r}^{\prime}_{j^{\prime}, S^{\prime}}} t_{j j^{\prime}} \left(\bm{r}_{j,S} - \bm{r}^{\prime}_{j^{\prime}, S^{\prime}}\right) e^{-i \bm{K}_{j, \xi} \cdot \bm{r}_{j,S}} e^{i \bm{K}_{j^{\prime, \xi^{\prime}}} \cdot \bm{r}^{\prime}_{j^{\prime, S^{\prime}}} } \hat{\psi}^{\dagger}_{j, S, \xi} \left(\bm{r}_{j,S}\right)  \hat{\psi}_{j^{\prime}, S^{\prime}, \xi^{\prime}} \left(\bm{r}_{j^{\prime}, S^{\prime}}\right).
\end{align}
这里我们将离散求和变成连续积分，Dirac comb为
\begin{align}
	\boxed{
		\sum_{\bm{r}_{j,S}} \delta^{\left(2\right)} \left(\bm{r} - \bm{r}_{j,S}\right) = \frac{1}{A_{\mathrm{mlg}}} \sum_{\bm{G}_{j}}e^{i \bm{G}_{j} \cdot \left(\bm{r} - \bm{\tau}_{j,S}\right)},
	}
\end{align}
这里\(\bm{G}_{j} = \mathcal{R} \left(\theta_{j}\right)\bm{G}\)是对应层的倒格矢。复式格子的Dirac comb 公式与格点的Dirac comb 公式略有不同，我们将在附录中给出证明。对于任意定义在第\(j\)层，\(S\)子晶格原子上的函数\(F\), 可以将求和\(\sum_{\bm{r}_{j,S}} F \left(\bm{r}_{j,S}\right)\)写成
\begin{align}
	\sum_{\bm{r}_{j,S}} F \left(\bm{r}_{j,S}\right) = \int d \bm{r} F \left(\bm{r}\right) \sum_{\bm{r}_{j,S}} \delta^{\left(2\right)} \left(\bm{r} - \bm{r}_{j,S}\right) = \frac{1}{A_{\mathrm{mlg}}} \sum_{\bm{G}_{j}} e^{-i \bm{G}_{j} \cdot \bm{\tau}_{j, S}} \int d \bm{r} e^{i \bm{G}_{j} \cdot \bm{r}} F \left(\bm{r}\right).
\end{align}
所以连续场算符可以表示成
\begin{align}
	\sum_{\bm{r}_{j,S}} \hat{\psi}_{j,S,\xi} \left(\bm{r}_{j,S}	\right) &\simeq \frac{1}{A_{\mathrm{mlg}}} \int d \bm{r} \sum_{\bm{G}_{j}} e^{i \bm{G}_{j} \cdot \left(\bm{r} - \bm{\tau}_{j,S}\right)} \hat{\psi}_{j,S,\xi} \left(\bm{r}\right) \\
	\sum_{\bm{r}^{\prime}_{j^{\prime}, S^{\prime}}} \hat{\psi}_{j^{\prime}, S^{\prime}, \xi^{\prime}} \left(\bm{r}^{\prime}_{j^{\prime}, S^{\prime}}\right) & \simeq  \frac{1}{A_{\mathrm{mlg}}} \int d \bm{r}^{\prime} \sum_{\bm{G}^{\prime}_{j^{\prime}}} e^{-i \bm{G}^{\prime}_{j^{\prime}} \cdot \left(\bm{r}^{\prime} - \bm{\tau}^{\prime}_{j^{\prime}, S^{\prime}}\right)} \hat{\psi}_{j^{\prime}, S^{\prime}, \xi^{\prime}} \left(\bm{r}^{\prime}\right).
\end{align}
这里我们采用了convention，也就是对于湮灭算符\(\hat{\psi}\), 我们设定其前面\(e\)指数上的系数是负的，也就是\(e^{-i \bm{G}^{\prime}_{j^{\prime}} \cdot \left(\bm{r}^{\prime} - \bm{\tau}^{\prime}\right)  }\)；对于产生算符\(\hat{\psi}^{\dagger}\)，我们设定其前面\(e\)指数上的系数是正的，也就是\(e^{i \bm{G}_{j} \cdot \left(\bm{r} - \bm{\tau}\right)}\)。之所以这里可以改变符号，是因为这里倒格矢关于反演是封闭的，也就是
\begin{align}
	\sum_{\bm{G}} e^{i \bm{G} \cdot \left( \bm{r} - \bm{\tau}\right)} = \sum_{\bm{G}} e^{-i \bm{G} \cdot \left( \bm{r} - \bm{\tau}\right)}.
\end{align}
习惯上对于产生湮灭算符取“一正一负”，是因为这与狄拉克符号的ket与bra的傅立叶变换是一致的，也就是\(e^{i \bm{k} \cdot \bm{r}} \cdots e^{i \bm{k}^{\prime} \cdot \bm{r}}\)。实际上，这里完全可以反过来进行标记，最后的物理结果并不会改变，仅仅相当于重新标记倒格矢，并且最后的转移动量保持一致。

由此，连续模型的哈密顿量是
\begin{align}
	\boxed{
		\begin{aligned}
			\begin{split}
				\hat{H}_{\mathrm{eff}} \simeq \frac{1}{A_{\mathrm{mlg}}} \sum_{j,j^{\prime}} \sum_{S,S^{\prime}} &\sum_{\xi,\xi^{\prime}} \sum_{\bm{G}_{j}, \bm{G}^{\prime}_{j^{\prime}}} \int d \bm{r} \int d \bm{r}^{\prime} e^{i \bm{G}_{j} \cdot \left(\bm{r} - \bm{\tau}_{j,S}\right)} e^{-i \bm{G}^{\prime}_{j^{\prime}} \cdot \left(\bm{r}^{\prime} - \bm{\tau}^{\prime}_{j^{\prime}, S^{\prime}}\right)} \\
				& t_{jj^{\prime}} \left(\bm{r} - \bm{r}^{\prime}\right) e^{-i \bm{K}_{j, \xi} \cdot \bm{r}} e^{i \bm{K}_{j^{\prime}, \xi^{\prime}} \cdot \bm{r}^{\prime}} \hat{\psi}^{\dagger}_{j,S,\xi} \left(\bm{r}\right) \hat{\psi}_{j^{\prime}, S^{\prime}, \xi^{\prime}} \left(\bm{r}^{\prime}\right).
			\end{split}
		\end{aligned}
	}
\end{align}

\subsection{证明不同谷之间的解耦}
将哈密顿量按照谷分块，
\begin{align}
	\hat{H}_{\mathrm{eff}} = \sum_{\xi,\xi^{\prime}} \hat{H}_{\xi \xi^{\prime}},
\end{align}
所以\(\xi = \xi^{\prime}\)对应intravalley的部分\(\hat{H}_{++}, \hat{H}_{--}\)，而\(\xi \neq \xi^{\prime}\)对应intravalley的部分\(\hat{H}_{+-}, \hat{H}_{-+}\)。接下来我们就要证明，在长波近似下，intravalley的部分是可以忽略的。为了分析积分中的空间振荡，定义两个原子位置之差
\begin{align}
	\bm{d} = \bm{r} - \bm{r}^{\prime},\quad \bm{r}^{\prime} = \bm{r} - \bm{d}.
\end{align}
将所有依赖\(\bm{r}, \bm{r}^{\prime}\)的指数相位整理起来
\begin{align}
	\begin{split}
		e^{i \bm{G}_{j} \cdot \bm{r}} e^{-i \bm{G}^{\prime}_{j^{\prime}} \cdot \bm{r}^{\prime}} e^{-i \bm{K}_{j, \xi} \cdot \bm{r}} e^{i \bm{K}_{j^{\prime}, \xi^{\prime}} \cdot \bm{r}^{\prime}} & = e^{i \bm{G}_{j} \cdot \bm{r}} e^{-i \bm{G}^{\prime}_{j^{\prime}} \cdot \left(\bm{r} - \bm{d}\right)}  e^{-i \bm{K}_{j, \xi} \cdot \bm{r}}  e^{i \bm{K}_{j^{\prime}, \xi^{\prime}} \cdot \left(\bm{r} - \bm{d}\right)} \\
		& = e^{i \bm{Q}_{\xi \xi^{\prime}}^{j j^{\prime}} \cdot \bm{r}} e^{i \left(\bm{G}^{\prime}_{j^{\prime}} - \bm{K}_{j^{\prime}, \xi^{\prime}} \right) \cdot \bm{d}},
	\end{split}
\end{align}
其中
\begin{align}
	\boxed{
		\bm{Q}_{\xi \xi^{\prime}}^{j j^{\prime}}  = \bm{G}_{j} - \bm{G}^{\prime}_{j^{\prime}} - \bm{K}_{j,\xi} + \bm{K}_{j^{\prime},\xi^{\prime}}. 
	}
\end{align}
将子晶格相位提出，哈密顿量为
\begin{align}
	\begin{split}
			\hat{H}_{\xi \xi^{\prime}} \simeq \frac{1}{A_{\mathrm{mlg}}} \sum_{j,j^{\prime}} \sum_{S,S^{\prime}} \sum_{\bm{G}_{j}, \bm{G}^{\prime}_{j^{\prime}}} e^{-i \bm{G}_{j} \cdot \bm{\tau}_{j,S}} e^{i \bm{G}^{\prime}_{j^{\prime}} \cdot \bm{\tau}^{\prime}_{j^{\prime}, S^{\prime}}} & \int d \bm{r}  \int d \bm{d} \, e^{i \bm{Q}_{\xi \xi^{\prime}}^{j j^{\prime}}  \cdot \bm{r}} e^{i \left(\bm{G}^{\prime}_{j^{\prime}} - \bm{K}_{j^{\prime}, \xi^{\prime}} \right) \cdot \bm{d}} \\
			& t_{jj^{\prime}} \left(\bm{d}\right) \hat{\psi}^{\dagger}_{j,S,\xi} \left(\bm{r}\right) \hat{\psi}_{j^{\prime}, S^{\prime}, \xi^{\prime}} \left(\bm{r} - \bm{d}\right).
	\end{split}
\end{align}
这里\(t \left(\bm{d}\right)\)随着\(\left|\bm{d}\right|\)快速衰减，同时\(\hat{\psi}^{\dagger} \left(\bm{r}\right) \hat{\psi} \left(\bm{r} - \bm{d}\right)\)随\(\bm{r}\)慢变。因此，决定哈密顿量振荡行为的就是\(e^{i \bm{Q}_{\xi \xi^{\prime}}^{j j^{\prime}}  \cdot \bm{r}}\). 如果\(\left|\bm{Q}\right| \sim a^{-1}\), 它在原子尺度快速振荡，连续积分中被抑制；只有\(\left|\bm{Q}\right| \ll a^{-1}\)，才可以保留。

对于层内的intravalley项，取\(j = j^{\prime},\, \xi = \xi^{\prime}\)，有
\begin{align}
	\bm{Q}_{\xi \xi^{\prime}}^{j j^{\prime}} = \bm{Q}_{\xi \xi}^{j j} = \bm{G}_{j} - \bm{G}^{\prime}_{j} - \bm{K}_{j,\xi} + \bm{K}_{j,\xi} = \bm{G}_{j} - \bm{G}^{\prime}_{j}
\end{align}
首先给出实空间中的分析，对于同一层的倒格矢，如果\(\bm{G}_{j} \neq \bm{G}^{\prime}_{j}\), 则
\begin{align}
	\left|\bm{G}_{j} - \bm{G}^{\prime}_{j}\right| \sim \left|\bm{b}_{1}\right| \sim \frac{1}{a}.
\end{align}
则被抑制，所以只有\(\bm{G}_{j} = \bm{G}^{\prime}_{j}\), 此时\(\left|\bm{Q}_{\xi \xi^{\prime}}^{j j^{\prime}}\right| = 0 \ll a^{-1}\)保留了下来。更直观的角度可以从动量守恒出发，对连续包络场做傅立叶变换
\begin{align}
	\boxed{
		\begin{aligned}
			\hat{\psi}^{\dagger}_{j,S,\xi} \left(\bm{r}\right) &= \int_{\left|\bm{k}\right| < \Lambda} \frac{d \bm{k}}{\left(2 \pi\right)^{2}} e^{-i \bm{k} \cdot \bm{r}} \hat{\psi}^{\dagger}_{j,S,\xi} \left(\bm{k}\right) , \\
			\hat{\psi}_{j^{\prime}, S^{\prime}, \xi^{\prime}} \left(\bm{r} - \bm{d}\right) & = \int_{\left|\bm{k}^{\prime}\right| < \Lambda} \frac{d \bm{k}^{\prime}}{\left(2 \pi\right)^{2}} e^{i \bm{k}^{\prime} \cdot \left(\bm{r} - \bm{d}\right)} \hat{\psi}_{j^{\prime}, S^{\prime}, \xi^{\prime}} \left(\bm{k}^{\prime}\right). 
		\end{aligned}
	}
\end{align}
中心坐标\(\bm{r} \)的完整e指数为
\begin{align}
	e^{i \left(\bm{Q}_{\xi \xi^{\prime}}^{j j^{\prime}} - \bm{k} + \bm{k}^{\prime}\right) \cdot \bm{r}}，
\end{align}
积分有
\begin{align}
	\int \mathrm{d} \bm{r} e^{i \left(\bm{Q}_{\xi \xi^{\prime}}^{j j^{\prime}} - \bm{k} + \bm{k}^{\prime}\right) \cdot \bm{r}} = \left(2 \pi\right)^{2} \delta \left(\bm{Q}_{\xi \xi^{\prime}}^{j j^{\prime}} - \bm{k} + \bm{k}^{\prime}\right).
\end{align}
由此得到动量匹配条件
\begin{align}
	\boxed{
		\bm{Q} = \bm{k} - \bm{k}^{\prime}.
	}
\end{align}
更一般的是
\begin{align}
	\boxed{
		\bm{G}_{j} - \bm{G}^{\prime}_{j^{\prime}} - \bm{K}_{j,\xi} + \bm{K}_{j^{\prime},\xi^{\prime}} = \bm{k}_{j} - \bm{k}_{j^{\prime}}.
	}
\end{align}
注意到
\begin{align}
	\left|\bm{k} - \bm{k}^{\prime}\right| < \left|\bm{k}\right| + \left|\bm{k}^{\prime}\right| < 2 \Lambda,
\end{align}
所以\(\left|\bm{Q}\right| < 2 \Lambda\)。而如果\(\bm{G}_{j} \neq \bm{G}^{\prime}_{j}\)，同样有
\begin{align}
	\left|\bm{Q}\right| = \left|\bm{G}_{j} - \bm{G}^{\prime}_{j}\right| \geq \frac{4 \pi}{\sqrt{3}a} > 2 \Lambda
\end{align}
所以只有
\begin{align}
	\bm{G}_{j} = \bm{G}^{\prime}_{j},
\end{align}
更进一步，
\begin{align}
	\bm{k} = \bm{k}^{\prime}.
\end{align}
对于层内的intervalley项，如果\(j^{\prime} = j, \quad \xi^{\prime} = - \xi\), 那么
\begin{align}
	\bm{Q}^{jj}_{\xi,-\xi} = \bm{G}_{j} - \bm{G}^{\prime}_{j} - \bm{K}_{j, \xi} + \bm{K}_{j, -\xi} = \bm{G}_{j} - \bm{G}^{\prime}_{j} - 2 \bm{K}_{j, \xi}.
\end{align}
长波条件要求
\begin{align}
	\bm{G}_{j} - \bm{G}^{\prime}_{j} \simeq 2 \bm{K}_{j, \xi},
\end{align}
但由于\(2 \bm{K}\)并不是倒格矢，而\(3 \bm{K}\)才是倒格矢，所以\(\left|\bm{Q}\right| \sim \frac{1}{a}\)。此项仍然在原子尺度快速振荡，层内的intervalley项被舍弃。

对于层间的intravalley项，取\(j \neq j^{\prime}, \quad \xi^{\prime} = \xi\), 有
\begin{align}
	\bm{Q}^{jj^{\prime}}_{\xi \xi} = \bm{G}_{j} - \bm{G}^{\prime}_{j^{\prime}} - \bm{K}_{j,\xi} + \bm{K}_{j^{\prime,\xi}},
\end{align}
由于小转角会带来一个相对小的动量差，所以不必一定要求\(\bm{G}_{j} = \bm{G}^{\prime}_{j^{\prime}}\)。同时有以下关系
\begin{align}
	\left|\bm{K}_{j^{\prime},\xi} - \bm{K}_{j,\xi}\right| \sim K \theta = \frac{4 \pi}{3 a} \theta \ll \frac{1}{a}.
\end{align}
所以可以选取倒格矢组合\(\left( \bm{G}_{j},\, \bm{G}^{\prime}_{j^{\prime}}\right)\)使得
\begin{align}
	\left|\bm{Q}^{jj^{\prime}}_{\xi \xi}\right| = \left|\bm{G}_{j} - \bm{G}^{\prime}_{j^{\prime}} - \bm{K}_{j,\xi} + \bm{K}_{j^{\prime,\xi}}\right| \sim K \theta ,
\end{align}
最低阶有三个由\(C_{3z}\)对称性联系的组合，所产生的三个小的剩余动量也就是后面的三个转移动量\(\bm{q}_{1},\,\bm{q}_{2},\, \bm{q}_{3}\)。(这里省略了谷指标)

对于层间的intervalley项，取\(j^{\prime} \neq j, \, \xi^{\prime} = -\xi \neq \xi \), 有
\begin{align}\
	\begin{split}
			\bm{Q}^{jj^{\prime}}_{\xi,-\xi} & = \bm{G}_{j} - \bm{G}^{\prime}_{j^{\prime}} - \bm{K}_{j,\xi} + \bm{K}_{j^{\prime, - \xi}} = \bm{G}_{j} - \bm{G}^{\prime}_{j^{\prime}} - \bm{K}_{j,\xi} - \bm{K}_{j^{\prime},\xi} \\
			& = \mathcal{R} \left(\theta_{j}\right) \left(\bm{G} - \bm{K}_{\xi}\right) - \mathcal{R}\left(\theta_{j^{\prime}}\right) \left(\bm{G}^{\prime} + \bm{K}_{\xi} \right),
	\end{split}
\end{align}
对于小转角有
\begin{align}
	\mathcal{\theta} \bm{G} = \bm{G} + \theta \hat{z} \times \bm{G} + O \left(\theta^{2} \left|\bm{G}\right|\right),
\end{align}
因此
\begin{align}
	\bm{Q}^{jj^{\prime}}_{\xi,-\xi} = \left(\bm{G} - \bm{G}^{\prime} - 2 \bm{K}_{\xi}\right) + \theta_{j} \hat{z} \times \left(\bm{G} - \bm{K}_{\xi}\right)	- \theta_{j^{\prime}} \hat{\bm{z}} \times \left(\bm{G}^{\prime} + \bm{K}_{\xi}\right) + O \left(\theta^{2}\right),
\end{align}
如果只考虑低阶倒格矢，如\(\left|\bm{G}\right|, \left|\bm{G}^{\prime}\right| \sim a^{-1}\),因此转角带来的动量差只有\(\left|\delta \bm{Q}\right| \sim \theta/a\), 而零转角时的动量失配为\(\left|\bm{Q}_{0}\right| = \left|\bm{G} - \bm{G}^{\prime} - 2 \bm{K}_{\xi}\right| \sim a^{-1}\)。在小角度条件下，\(\theta \ll 1\), 因此，低阶倒格矢的旋转修正，不可能抵消原本\(O\left(a^{-1}\right)\)的动量失配。

但是如果选取非常高阶的倒格矢，使得\(\left|\bm{G}\right|, \left|\bm{G}^{\prime}\right| \sim 1/a \theta\)。那么旋转造成的偏移即可达到\(\theta \left|\bm{G}\right| \sim 1/a\), 这就可以补偿原本\(O \left(1/a\right)\)的动量失配。因此在高阶情况下，原则上可能存在倒格矢对使得
\begin{align}
	\left|\mathcal{R}\left(\theta_{j}\right) \bm{G} - \mathcal{R} \left(\theta_{j^{\prime}}\right) \bm{G}^{\prime} - \bm{K}_{j,\xi} - \bm{K}_{j^{\prime},\xi}\right| < 2 \Lambda,
\end{align}
但是对于高阶情况下，\(t \left(\bm{d}\right)\) 傅里叶变换后的\(\tilde{t }\left(\bm{p }\right)\)随着\(\left|\bm{p }\right|\)的增大迅速衰减。因此，层间intervalley项在连续模型中可以忽略。谷解耦
\begin{align}
	\hat{H }_{\mathrm{eff}} \simeq \hat{H }_{+} \oplus \hat{H }_{-}.
\end{align}
两个谷由时间反演操作\(\mathcal{T }\)联系。因此，只需要固定一个谷， 我们一般选取\(\xi = +\)。固定谷后，哈密顿量可以写成
\begin{align}
	\begin{split}
		\boxed{
			\begin{aligned}
				\hat{H}_{\xi \xi^{\prime}} \simeq \frac{1}{A_{\mathrm{mlg}}} \sum_{j,j^{\prime}} \sum_{S,S^{\prime}} \sum_{\bm{G}_{j}, \bm{G}^{\prime}_{j^{\prime}}} e^{-i \bm{G}_{j} \cdot \bm{\tau}_{j,S}} e^{i \bm{G}^{\prime}_{j^{\prime}} \cdot \bm{\tau}^{\prime}_{j^{\prime}, S^{\prime}}} & \int \mathrm{d} \bm{r} \int \mathrm{d} \bm{d} \, e^{i \bm{Q}_{\xi \xi^{\prime}}^{j j^{\prime}}  \cdot \bm{r}} e^{i \left(\bm{G}^{\prime}_{j^{\prime}} - \bm{K}_{j^{\prime}, \xi} \right) \cdot \bm{d}} \\
				& t_{jj^{\prime}} \left(\bm{d}\right) \hat{\psi}^{\dagger}_{j,S,\xi} \left(\bm{r}\right) \hat{\psi}_{j^{\prime}, S^{\prime}, \xi} \left(\bm{r} - \bm{d}\right).
			\end{aligned}
		}
	\end{split}
\end{align}

\subsection{梯度展开}
接下来也就是对\(\hat{\psi}\left(\bm{r} - \bm{d }\right)\)做梯度展开，展开到一阶
\begin{align}
	\hat{\psi}_{j^{\prime}, S^{\prime}, \xi} \left(\bm{r} - \bm{d}\right) = e^{-\bm{d } \cdot \nabla_{\bm{r}}} \hat{\psi}_{j^{\prime}, S^{\prime}, \xi} \left(\bm{r}\right) \simeq  \hat{\psi}_{j^{\prime}, S^{\prime}, \xi} \left(\bm{r}\right)  - d_{\mu} \partial_{\mu}  \hat{\psi}_{j^{\prime}, S^{\prime}, \xi} \left(\bm{r}\right),
\end{align}
（这里当然也可以像Kang的文章中用对称坐标，对\(\psi^{\dagger} \left(\bm{r} + \bm{d}/2\right) \psi \left(\bm{r} - \bm{d}/2\right)\)）进行展开。首先对于层内项(interlayer coupling)\(j = j^{\prime}\), 有\(\bm{G}_{j} = \bm{G}^{\prime}_{j}\)，以及\(\bm{Q } = 0\), 哈密顿量可以写为
\begin{align}
	\hat{H}^{\mathrm{intra}}_{j,\xi} \simeq \frac{1}{A_\mathrm{mlg}} \sum_{S,S^{\prime}} \sum_{\bm{G}_{j}} e^{-i \bm{G}_{j} \cdot \left( \bm{\tau}_{j,S } - \bm{\tau}^{\prime}_{j, S^{\prime}}\right)} \int \mathrm{d} \bm{r} \mathrm{d} \bm{r} e^{i \left(\bm{G}_{j} - \bm{K }_{j,\xi}\right) \cdot \bm{d }} \, t_{jj} \left(\bm{d }\right) \hat{\psi}^{\dagger}_{j, S, \xi} \left(\bm{r}\right) \hat{\psi}_{j,S^{\prime}, \xi} \left(\bm{r} - \bm{d }\right).
\end{align}
将连续场展开到一阶，
\begin{align}
	\hat{H}^{\mathrm{intra}}_{j,\xi} \simeq \int \mathrm{d} \bm{r} \sum_{S,S^{\prime}} \hat{\psi}^{\dagger}_{j, S, \xi} \left(\bm{r}\right) \left[\mathcal{M }^{\left(0\right)}_{j,\xi;SS^{\prime}} - \mathcal{M }^{\left(1\right)}_{j,\xi;SS^{\prime}, \mu} \partial_{\mu}\right] \hat{\psi}_{j,S^{\prime},\xi} \left(\bm{r}\right),
\end{align}
其中
\begin{align}
	\mathcal{M }^{\left(0\right)}_{j,\xi;SS^{\prime}} = \frac{1}{A_{\mathrm{mlg}}} \sum_{\bm{G}_{j}}  e^{-i \bm{G}_{j} \cdot \left( \bm{\tau}_{j,S } - \bm{\tau}^{\prime}_{j, S^{\prime}}\right)} \int \mathrm{d} \bm{d} \, t_{jj} \left(\bm{d }\right) e^{i \left(\bm{G}_{j} - \bm{K }_{j,\xi}\right) \cdot \bm{d }},
\end{align}
以及
\begin{align}
	\mathcal{M }^{\left(1\right)}_{j,\xi;SS^{\prime}, \mu}  = \frac{1}{A_{\mathrm{mlg}}} \sum_{\bm{G}_{j}}  e^{-i \bm{G}_{j} \cdot \left( \bm{\tau}_{j,S } - \bm{\tau}^{\prime}_{j, S^{\prime}}\right)} \int \mathrm{d} \bm{d} \, d_{\mu} t_{jj} \left(\bm{d }\right) e^{i \left(\bm{G}_{j} - \bm{K }_{j,\xi}\right) \cdot \bm{d }}.
\end{align}
可以看到\(\mathcal{M }^{\left(0\right)}_{j,\xi}\)和\(\mathcal{M }^{\left(1\right)}_{j,\xi;\mu}\)都是作用在\(A,B\)子晶格空间的\(2 \times 2\)矩阵。我们这里需要强调梯度展开本身给出的并不是\(\bm{\sigma} \cdot \hat{\bm{p}},\quad \hat{\bm{p}} = - i \nabla_{\bm{r}}\), 梯度展开给出的一阶项必然具有以下形式
\begin{align}
	\hat{H}^{\mathrm{intra}, \left(1\right)}_{j,\xi} = -\\int \mathrm{d} \bm{r} \hat{\psi}^{\dagger}_{j, \xi} \left(\bm{r}\right) \left(\mathcal{M }^{\left(1\right)}_{j,\xi;x} \partial_{x } + \mathcal{M }^{\left(1\right)}_{j,\xi;y} \partial_{y}\right) \hat{\psi}_{j,\xi} \left(\bm{r}\right)
\end{align}
写成更一般的形式也就是
\begin{align}
	h^{\left(1\right)}_{j,\xi} = M_x \left(- i \partial_{x } \right) + M_{y} \left(-i \partial_{y}\right).
\end{align}
其中\(M_{x }, M_{y}\)是厄米矩阵。那么接下来也就是通过对称性来确定\(\mathcal{M}^{\left(0\right)}\)和\(\mathcal{M }^{\left(1\right)}\)的形式。\footnote{实在写不动了，以后有机会再写吧，2026年7月29日}

定义子晶格的两分量场，
\begin{align}
	\hat{\psi}_{j,\xi} \left(\bm{r}\right) = \begin{pmatrix}
		\hat{\psi}_{j,A,\xi} \left(\bm{r}\right) \\ \hat{\psi}_{j,B,\xi} \left(\bm{r}\right)
	\end{pmatrix}
\end{align}
层内的零阶项可以被泡利矩阵分解成
\begin{align}
	\mathcal{M }^{\left(0\right)}_{j,\xi} = a_0 \sigma_{0} + a_{1} \sigma_{1} + a_2 \sigma_{2} + a_3 \sigma_{3}.
\end{align}
对于仅考虑最近邻跃迁的单层石墨烯紧束缚模型, 哈密顿量只包含连接 A/B 两个子晶格的跃迁项。由于不存在同一子晶格内的跃迁，哈密顿量不含与 \(\sigma_{0}\) 成正比的动量依赖项；同时可将狄拉克点能量选为能量零点，从而略去常数形式的 \(\sigma_{0}\) 项。同时，由于A/B 子晶格等价，二者具有相同的在位能，因此也不存在\(\sigma_{3}\)质量项。另外，在谷中心\(\bm{K }_{\xi}\)处，三个由\(C_{3z}\)旋转互相关联的最近邻跃迁具有相同的跃迁振幅，而相应的相位依次相差\(e^{i 2 \pi/3}\)，从而有
\begin{align}
	1 + e^{i 2 \pi/3} + e^{ - i 2 \pi/3} = 0,
\end{align}
所以谷中心处哈密顿量的非对角矩阵元同样为零。综上
\begin{align}
	\mathcal{M }^{\left(0\right)}_{j,\xi} = 0.
\end{align}

对于一阶项\(\mathcal{M }^{\left(1\right)}\),

\subsection{平面波表示}


\section{附录}
\subsection{其他体系中的紧束缚模型}
如果是复式晶格，则格点的费米子算符就需要加入子晶格自由度，也就是
\begin{align}
	\hat{H} = \sum_{\bm{R}, \bm{\delta}} t_{s,s^{\prime}} \left(\bm{\delta}\right) \hat{c}^{\dagger}_{\bm{R}, s} \hat{c}_{\bm{R} - \bm{\delta}, s^{\prime}} + h.c.
\end{align}
这里我们通常会假定跃迁只与格点之间的相对位置有关，虽然我们这里加入了子晶格自由度\(s\)，但是\(\hat{c}^{\dagger}_{\bm{R}, s} \) 还是格点的费米子算符，只不过是\(s\) 子晶格上的格点的费米子算符。

同时我们这里并没有像Kang的文章中对于晶格形变所讨论的那样，由于晶格形变改变了跃迁所依赖的实际原子位置，而不仅仅是格点位置，所以我们需要用原子费米子算符来写出紧束缚模型的哈密顿量，也就是
\begin{align}
	\hat{H} = \sum_{s s^{\prime}}\sum_{\bm{r}_{s}，\bm{r}^{\prime}_{s^{\prime}}} t\left(\bm{X}_{s} - \bm{X}^{\prime}_{s^{\prime}}\right) \hat{c}^{\dagger}_{s, \bm{r}_{s}}  \hat{c}_{s^{\prime}, \bm{r}^{\prime}_{s^{\prime}}}.
\end{align}
这里\(\bm{X}_{s}, \bm{X}^{\prime}_{s^{\prime}}\)是形变后的原子坐标，这里不多加赘述。


\subsection{何为有效模型？}
有效模型中的有效，指的是只保留所关心的低能自由度，忽略或积分掉高能自由度，从而近似描述某个有限能量和动量范围内的物理。

假设完整的哈密顿量是\(\hat{H}\), 定义低能子空间的投影算符是\(\mathcal{P}\), 高能子空间的投影算符为\(\mathcal{Q} = 1 - 	\mathcal{P}\)。	有效哈密顿量只作用在低能子空间当中，最简单的投影形式是
\begin{align}
	\hat{H}_{\mathrm{eff}} \simeq \mathcal{P} \hat{H} \mathcal{P},
\end{align}
如果还考虑高能态对低能态的虚跃迁，则更一般地还会出现
\begin{align}
	\hat{H}_{\mathrm{eff}} \left(E\right) = \mathcal{P} \hat{H} \mathcal{P} + \mathcal{P} \hat{H} \mathcal{Q} \frac{1}{E - \mathcal{Q} \hat{H} \mathcal{Q}} \mathcal{Q} \hat{H} \mathcal{P} + \cdots,
\end{align}
因此，高能自由度虽然不再被显式求解，但是它们的影响可能被吸收到有效速度，有效质量，耦合常数等参数当中。因此，有效意味着Hilbert空间被缩小了，只保留低能自由度；同时哈密顿量只保留在小参数下最重要的有限阶项。

对于\(\bm{k} \cdot \bm{p}\)模型。严格历史意义上的\(\bm{k} \cdot \bm{p}\)名称来源于连续晶体哈密顿量中的
\begin{align}
	\frac{\left(\hat{\bm{p}}+ \hbar \bm{k}\right)^{2}}{2 m} = \frac{\hat{\bm{p}}^{2}}{2m} + \frac{\hbar}{m} \bm{k} \cdot \hat{\bm{p}} + \frac{\hbar^{2} k^{2}} {2m},
\end{align}
这里出现了真正的\(\bm{k} \cdot \bm{p}\)项。但是，现在\(\bm{k} \cdot \bm{p}\)模型更宽泛地指在某个高对称动量点附近构造的低能哈密顿量矩阵。比如，假设只关心\(\bm{K}\)附近的状态，写成\(\bm{k} = \bm{K} + \bm{q}, \, \left|\bm{q}\right| a \ll 1\). 对布洛赫哈密顿量做长波展开有
\begin{align}
	h \left(\bm{K} + \bm{q}\right) = h \left(\bm{K}\right) + q_i \frac{\partial h}{\partial k_i} \bigg|_{\bm{k} = \bm{K}} + \frac{1}{2} q_i q_j \frac{\partial^{2} h}{\partial k_i \partial k_j} \bigg|_{\bm{k} = \bm{K}} ,
\end{align}
然后通常只保留有限阶项，得到有效布洛赫哈密顿量
\begin{align}
	h_{\mathrm{eff}} \left(\bm{q}\right) = h \left(\bm{K}\right) + q_i \frac{\partial h}{\partial k_i} \bigg|_{\bm{k} = \bm{K}}.
\end{align}
所以有效模型及意味着保留低能部分。

\subsection{Possion 求和公式（Dirac comb forumla）的证明}
考虑周期系统中的密度函数, 对于简单晶格的情况
\begin{align}
	\rho \left(\bm{r}\right) = \sum_{\bm{R}} \delta^{\left(2\right)}  \left(\bm{r} - \bm{R}\right).
\end{align}
这是一个与晶格具有相同周期性的函数，因此可以按倒格矢做傅立叶展开
\begin{align}
	\rho \left(\bm{r}\right) = \sum_{\bm{G}} c_{\bm{G}} e^{i \bm{G}\cdot\bm{r}}.
\end{align}
傅立叶系数为
\begin{align}
	c_{\bm{G}} = \frac{1}{A_{\mathrm{mlg}}} \int_{\mathrm{uc}} \rho \left(\bm{r}\right) e^{-i \bm{G} \cdot \bm{r}} = \frac{1}{A_{\mathrm{mlg}}},
\end{align}
因为一个原胞内只贡献一个delta峰。所以有Possion 求和公式
\begin{align}
	\boxed{
		\sum_{\bm{R}} \delta^{\left(2\right)} \left(\bm{r} - \bm{R}\right)	= \frac{1}{A_{\mathrm{mlg}}} \sum_{\bm{G}} e^{i \bm{G} \cdot \bm{r}}.
	}
\end{align}
对于复式晶格，原胞内子晶格\(S\)的位置为\(\bm{\tau}_{S}\), 密度函数需要改写成
\begin{align}
	\rho_{S} \left(\bm{r}\right) = \sum_{\bm{R}} \delta^{\left(2\right)} \left(\bm{r} - \bm{R} - \bm{\tau}_{S}\right) = \rho \left(\bm{r} - \bm{\tau}_{S}\right),
\end{align}
将上式中的\(\bm{r}\)替换成\(\bm{r} - \bm{\tau}_{S}\), 将\(\bm{R}\)替换成实际原子位置\(\bm{r}_{S}\)。得到
\begin{align}
	\boxed{
		\sum_{\bm{r}_{S}} \delta^{\left(2\right)} \left(\bm{r} - \bm{r}_{S}\right) = \frac{1}{A_{\mathrm{mlg}}} \sum_{\bm{G}} e^{i \bm{G} \cdot \left(\bm{r} - \bm{\tau}_{S}\right)}.
	}
\end{align}
即子晶格的平移等价于傅立叶展开中新增加的相位。
\end{document}
